% !TEX root = ../main.tex
\begin{abstract}[zh]
\addcontentsline{toc}{chapter}{摘 \quad 要}
现代 3C 电子装配线需要处理小型元器件、高产品多样性和富接触操作，这些任务仍难以通过传统固定程序机器人实现自动化。本项目针对这一缺口，面向双臂工业操作开发视觉-语言-动作（Vision-Language-Action, VLA）系统。系统设计运行于双臂 Franka 平台，并面向典型 PC 关键部件任务，包括内存条插入、CPU 放置以及电源线连接。
\par
当前部署系统选择相应的 ACT 检查点，接收同步视觉观测和机器人状态，并在闭环执行过程中输出机器人及夹爪动作。现阶段证据限于干运行和有边界的验证活动；本文不声称已经实现语言条件推理、统一多任务控制、自主规划、故障检测、恢复或双臂协同操作。这些功能仍是面向预期 VLA 与双臂系统的长期扩展方向。
\par
工程验收目标包括：规划成功率高于 90\%，任务执行成功率高于 70\%，部件检测精度高于 90\%，位姿误差低于 20\,mm，故障检测率高于 90\%，恢复成功率高于 70\%。这些数值定义验证阈值，而非已经证明的性能结果。
\end{abstract}

{
  \newfontfamily{\arial}{Arial}[
    Scale=0.94,
    BoldFont={Arial Bold},
    ItalicFont={Arial Italic},
    BoldItalicFont={Arial Bold Italic},
  ]
\ctexset{chapter/format+={\arial}}
\begin{abstract}[en]
\addcontentsline{toc}{chapter}{ABSTRACT}
Modern 3C electronics assembly lines deal with small components, high product variety, and contact-rich operations that remain difficult to automate with conventional fixed-program robots. Our project targets this gap by developing a Vision-Language-Action (VLA) system for bimanual industrial manipulation. The system is designed to run on a dual-arm Franka platform and to handle representative PC key-component tasks---specifically RAM insertion, CPU placement, and power-cable connection.
\par
The deployed system selects the corresponding ACT checkpoint, receives synchronized visual observations and robot state, and outputs robot and gripper actions during closed-loop execution. Current evidence is limited to dry runs and bounded validation activities; the report does not claim demonstrated language-conditioned inference, unified multi-task control, autonomous planning, failure detection, recovery, or bimanual manipulation. These functions remain long-term extensions toward the intended Vision-Language-Action (VLA) and dual-arm system.
\par
The engineering acceptance targets are planning success above 90\%, task execution success above 70\%, component detection accuracy above 90\%, pose error below 20\,mm, failure-detection rate above 90\%, and recovery success above 70\%. These values define validation thresholds rather than demonstrated performance results.
\end{abstract}
}
